% Sample documents: running text in many languages, sizes and features, to
% judge colour, spacing and the drafted glyphs in use.
%   lualatex --output-directory=build/proofs proofs/samples.tex   (or make proof)
\documentclass[11pt]{article}
\usepackage[a4paper,margin=2.4cm]{geometry}
\usepackage{fontspec}
\setmainfont{MyiaDidot-Regular.otf}[Path=../fonts/otf/, Renderer=HarfBuzz,
  BoldFont=MyiaDidot-Bold.otf, ItalicFont=MyiaDidot-Italic.otf, BoldItalicFont=MyiaDidot-BoldItalic.otf]
\newfontfamily\osf{MyiaDidot-Regular.otf}[Path=../fonts/otf/, Renderer=HarfBuzz, Numbers=OldStyle,
  BoldFont=MyiaDidot-Bold.otf, ItalicFont=MyiaDidot-Italic.otf, BoldItalicFont=MyiaDidot-BoldItalic.otf]
\newcommand\smc[1]{{\addfontfeature{Letters=SmallCaps}#1}}
\newcommand\allsc[1]{{\addfontfeature{Letters={SmallCaps,UppercaseSmallCaps}}#1}}
\newcommand\lang[2]{{\addfontfeature{Language=#1}#2}}
\newfontlanguage{IPA}{IPPH}
\newcommand\head[1]{\par\bigskip{\large\allsc{#1}}\par\smallskip}
\newcommand\bi[1]{{\bfseries\itshape #1}\par}
\newcommand\lng[1]{{\addfontfeature{Letters={SmallCaps,UppercaseSmallCaps}}#1.}\enspace}
\setlength{\parindent}{1.2em}
\linespread{1.08}
\pagestyle{empty}
\begin{document}

% -- 1. a book page ---------------------------------------------------------
\begin{center}
{\Huge Pride and Prejudice}\par\medskip
{\large\smc{by Jane Austen}}\par\bigskip
{\allsc{Chapter I}}
\end{center}
\begingroup\osf
It is a truth universally acknowledged, that a single man in possession of a
good fortune, must be in want of a wife.

However little known the feelings or views of such a man may be on his first
entering a neighbourhood, this truth is so well fixed in the minds of the
surrounding families, that he is considered the rightful property of some one
or other of their daughters.

\smc{“My dear Mr.\ Bennet,”} said his lady to him one day, “have you heard that
Netherfield Park is let at last?”

Mr.\ Bennet replied that he had not. “But it is,” returned she; “for Mrs.\ Long
has just been here, and she told me all about it.” Mr.\ Bennet made no answer.
“Do you not want to know who has taken it?” cried his wife impatiently.
“\emph{You} want to tell me, and I have no objection to hearing it.” This was
invitation enough.\footnote{First published in 1813 by Thomas Egerton,
London, in three volumes; the text here follows the first edition, p.~1--2.}

\bigskip
\begin{center}{\allsc{Du côté de chez Swann}}\par\smallskip{\smc{Marcel Proust, 1913}}\end{center}
\lang{French}{Longtemps, je me suis couché de bonne heure. Parfois, à peine ma
bougie éteinte, mes yeux se fermaient si vite que je n’avais pas le temps de me
dire : « Je m’endors. » Et, une demi-heure après, la pensée qu’il était temps de
chercher le sommeil m’éveillait ; je voulais poser le volume que je croyais
avoir encore dans les mains et souffler ma lumière ; je n’avais pas cessé en
dormant de faire des réflexions sur ce que je venais de lire, mais ces
réflexions avaient pris un tour un peu particulier ; il me semblait que
j’étais moi-même ce dont parlait l’ouvrage : une église, un quatuor, la
rivalité de François~I\textsuperscript{er} et de Charles-Quint.}
\endgroup

\newpage
% -- 2. one text, many languages -------------------------------------------
\begin{center}{\large\allsc{Universal Declaration of Human Rights, Article 1}}\end{center}
\setlength{\parindent}{0pt}\setlength{\parskip}{5pt}

\lng{English} All human beings are born free and equal in dignity and
rights. They are endowed with reason and conscience and should act towards one
another in a spirit of brotherhood.

\lng{Français} \lang{French}{Tous les êtres humains naissent libres et égaux
en dignité et en droits. Ils sont doués de raison et de conscience et doivent
agir les uns envers les autres dans un esprit de fraternité.}

\lng{Deutsch} \lang{German}{Alle Menschen sind frei und gleich an Würde und
Rechten geboren. Sie sind mit Vernunft und Gewissen begabt und sollen einander
im Geist der Brüderlichkeit begegnen.}

\lng{Polski} Wszyscy ludzie rodzą się wolni i równi pod względem swej
godności i swych praw. Są oni obdarzeni rozumem i sumieniem i powinni
postępować wobec innych w duchu braterstwa.

\lng{Čeština} Všichni lidé rodí se svobodní a sobě rovní co do důstojnosti a
práv. Jsou nadáni rozumem a svědomím a mají spolu jednat v duchu bratrství.

\lng{Magyar} Minden emberi lény szabadon születik és egyenlő méltósága és
joga van. Az emberek, ésszel és lelkiismerettel bírván, egymással szemben
testvéri szellemben kell hogy viseltessenek.

\lng{Română} \lang{Romanian}{Toate ființele umane se nasc libere și egale în
demnitate și în drepturi. Ele sunt înzestrate cu rațiune și conștiință și
trebuie să se comporte unele față de altele în spiritul fraternității.}

\lng{Lietuvių} Visi žmonės gimsta laisvi ir lygūs savo orumu ir teisėmis.
Jiems suteiktas protas ir sąžinė ir jie turi elgtis vienas kito atžvilgiu kaip
broliai.

\lng{Íslenska} Hver maður er borinn frjáls og jafn öðrum að virðingu og
réttindum. Menn eru gæddir vitsmunum og samvisku, og ber þeim að breyta
bróðurlega hverjum við annan.

\lng{Türkçe} \lang{Turkish}{Bütün insanlar hür, haysiyet ve haklar
bakımından eşit doğarlar. Akıl ve vicdana sahiptirler ve birbirlerine karşı
kardeşlik zihniyeti ile hareket etmelidirler.}

\lng{Azərbaycanca} \lang{Azeri}{Bütün insanlar ləyaqət və hüquqlarına görə
azad və bərabər doğulurlar. Onların şüurları və vicdanları var və
bir-birlərinə münasibətdə qardaşlıq ruhunda davranmalıdırlar.}

\lng{Tiếng Việt} \lang{Vietnamese}{Tất cả mọi người sinh ra đều được tự do và
bình đẳng về nhân phẩm và quyền. Mọi con người đều được tạo hóa ban cho lý trí
và lương tâm và cần phải đối xử với nhau trong tình bằng hữu.}

\lng{Ελληνικά} Όλοι οι άνθρωποι γεννιούνται ελεύθεροι και ίσοι στην
αξιοπρέπεια και τα δικαιώματα. Είναι προικισμένοι με λογική και συνείδηση, και
οφείλουν να συμπεριφέρονται μεταξύ τους με πνεύμα αδελφοσύνης.

\lng{Русский} Все люди рождаются свободными и равными в своем достоинстве и
правах. Они наделены разумом и совестью и должны поступать в отношении друг
друга в духе братства.

\lng{Українська} Всі люди народжуються вільними і рівними у своїй гідності
та правах. Вони наділені розумом і совістю і повинні діяти у відношенні один
до одного в дусі братерства.

\lng{Български} Всички хора се раждат свободни и равни по достойнство и
права. Те са надарени с разум и съвест и следва да се отнасят помежду си в дух
на братство.

\lng{Српски} Сва људска бића рађају се слободна и једнака у достојанству и
правима. Она су обдарена разумом и свешћу и треба једни према другима да
поступају у духу братства.

\smallskip
\bi{Čeština: Všichni lidé rodí se svobodní · Magyar: Minden emberi lény · Polski: Wszyscy ludzie}
\bi{\lang{Romanian}{Română: Toate ființele umane se nasc libere} · Lietuvių: Visi žmonės gimsta laisvi}
\bi{\lang{Turkish}{Türkçe: Bütün insanlar hür} · \lang{Azeri}{Azərbaycanca: Bütün insanlar azad} · Íslenska: Hver maður}
\bi{\lang{Vietnamese}{Tiếng Việt: Tất cả mọi người sinh ra đều được tự do và bình đẳng}}
\bi{Ελληνικά: Όλοι οι άνθρωποι γεννιούνται ελεύθεροι · \lang{Serbian}{Српски: Сва људска бића}}

\newpage
% -- 3. Vietnamese at length ---------------------------------------------------
\begin{center}{\large\allsc{Truyện Kiều}}\par\smallskip\smc{Nguyễn Du}\end{center}
\lang{Vietnamese}{%
\begin{center}
Trăm năm trong cõi người ta,\\
Chữ tài chữ mệnh khéo là ghét nhau.\\
Trải qua một cuộc bể dâu,\\
Những điều trông thấy mà đau đớn lòng.
\end{center}
\setlength{\parindent}{1.2em}
{\Large ẤẦẨẪẬ ẮẰẲẴẶ ẾỀỂỄỆ ỐỒỔỖỘ ỚỜỞỠỢ ỨỪỬỮỰ}\par
{\Large ấầẩẫậ ắằẳẵặ ếềểễệ ốồổỗộ ớờởỡợ ứừửữự ỳýỷỹỵ đ Đ}\par
\allsc{Tiếng Việt: Người ấy ở Huế, đẹp như tranh vẽ.}\par
\bi{Trăm năm trong cõi người ta, chữ tài chữ mệnh khéo là ghét nhau.}
\bi{{\Large ẤẦẨẪẬ ẮẰẲẴẶ ẾỀỂỄỆ ỐỒỔỖỘ ỚỜỞỠỢ ỨỪỬỮỰ}}
\bi{{\Large ấầẩẫậ ắằẳẵặ ếềểễệ ốồổỗộ ớờởỡợ ứừửữự ỳýỷỹỵ đ Đ}}}

% -- 4. African orthographies, IPA ----------------------------------------------
\head{African orthographies}
\lng{Hausa} ƙasa, ɗaki, ɓera, ƴaƴa — Ƙ Ɗ Ɓ Ƴ\quad
\lng{Ewe} Eʋegbe, ɖevi, ƒe, ɛ, ɔ — Ʋ Ɖ Ƒ Ɛ Ɔ\quad
\lng{Yorùbá} ẹ̀kọ́, ọmọ, ṣé — Ẹ Ọ Ṣ\quad
\lng{Lingála} mbɔ́tɛ\quad
\lng{Fulfulde} ɓe, ɗum, ŋ, ƴ — Ɓ Ɗ Ŋ Ƴ\quad
\lng{Kabiyè and others} Ɖ ɖ Ɛ ɛ Ɣ ɣ Ɩ ɩ Ŋ ŋ Ɔ ɔ Ʊ ʊ Ʒ ʒ Ɑ ɑ Ʋ ʋ\par
\bi{Hausa ƙasa, ɗaki, ɓera, ƴaƴa — Ƙ Ɗ Ɓ Ƴ\quad Ewe Eʋegbe, ɖevi, ƒe, ɛ, ɔ — Ʋ Ɖ Ƒ Ɛ Ɔ}
\bi{Yorùbá ẹ̀kọ́, ọmọ, ṣé — Ẹ Ọ Ṣ\quad Fulfulde ɓe, ɗum, ŋ, ƴ — Ɓ Ɗ Ŋ Ƴ\quad Ɣ ɣ Ɩ ɩ Ʊ ʊ Ʒ ʒ Ɑ ɑ}

\head{The North Wind and the Sun, in IPA}
\begingroup\addfontfeature{Language=IPA}% the two-storey a in the italics too
\def\aaa{a ɑ α · sat sɑt · ɐ æ ʌ · θ ϑ · ɪ i · ʏ y · ɴ n · ʀ r · ɢ g · ʟ l · ʙ b · ʜ h}
\aaa\par{\itshape\aaa\par}{\bfseries\aaa\par}{\bfseries\itshape\aaa\par}\smallskip
ðə ˈnɔɹθ ˌwɪnd ən ðə ˈsʌn wɚ dɪˈspjutɪŋ ˈwɪtʃ wəz ðə ˈstɹɔŋɡɚ, wɛn ə
ˈtɹævlɚ ˌkeɪm əˈlɑŋ ˈɹæpt ɪn ə ˈwɔɹm ˈkloʊk. ðeɪ əˈɡɹiːd ðət ðə ˈwʌn hu
ˈfɚst səkˈsidəd ɪn ˈmeɪkɪŋ ðə ˈtɹævlɚ ˈteɪk hɪz ˈkloʊk ˌɔf ʃʊd bi
kənˈsɪdɚd ˈstɹɔŋɡɚ ðən ðə ˈʌðɚ.

\smallskip
ɓ ɗ ʄ ɠ ʛ · ɖ ɭ ɳ ɽ ʈ ʂ ʐ ɻ · ʋ ɹ ɰ ɥ ɯ · ʃ ʒ ɕ ʑ ɧ · ɸ β θ ð χ ʁ ħ ʕ ʔ ʡ ʢ ·
ɚ ɝ ɞ ɘ ɵ ɤ ɶ ʉ ɨ · ǀ ǁ ǂ ǃ ʘ · tʰ tʲ tʷ nˠ aː aˑ ə̃ n̩ t̪ d̥ e̞ o̝ ɑ̟ ʊ̈ ˥˦˧˨˩\par
{\itshape ðə ˈnɔɹθ ˌwɪnd ən ðə ˈsʌn wɚ dɪˈspjutɪŋ ˈwɪtʃ wəz ðə ˈstɹɔŋɡɚ, wɛn ə
ˈtɹævlɚ ˌkeɪm əˈlɑŋ ˈɹæpt ɪn ə ˈwɔɹm ˈkloʊk. ðeɪ əˈɡɹiːd ðət ðə ˈwʌn hu
ˈfɚst səkˈsidəd ɪn ˈmeɪkɪŋ ðə ˈtɹævlɚ ˈteɪk hɪz ˈkloʊk ˌɔf ʃʊd bi
kənˈsɪdɚd ˈstɹɔŋɡɚ ðən ðə ˈʌðɚ.\par
ɓ ɗ ʄ ɠ ʛ · ɖ ɭ ɳ ɽ ʈ ʂ ʐ ɻ · ʋ ɹ ɰ ɥ ɯ · ʃ ʒ ɕ ʑ ɧ · ɸ β θ ð χ ʁ ħ ʕ ʔ ʡ ʢ\par
ɚ ɝ ɞ ɘ ɵ ɤ ɶ ʉ ɨ · ǀ ǁ ǂ ǃ ʘ · tʰ tʲ tʷ nˠ aː aˑ ə̃ n̩ t̪ d̥ e̞ o̝ ɑ̟ ʊ̈ ˥˦˧˨˩\par}
\bi{ðə ˈnɔɹθ ˌwɪnd ən ðə ˈsʌn wɚ dɪˈspjutɪŋ ˈwɪtʃ wəz ðə ˈstɹɔŋɡɚ}
\bi{ɓ ɗ ʄ ɠ ʛ · ɖ ɭ ɳ ɽ ʈ ʂ ʐ ɻ · ʋ ɹ ɰ ɥ ɯ · ʃ ʒ ɕ ʑ ɧ · ɸ β θ ð χ ʁ ħ ʕ ʔ ʡ ʢ}
\bi{ɚ ɝ ɞ ɘ ɵ ɤ ɶ ʉ ɨ · ǀ ǁ ǂ ǃ ʘ · tʰ tʲ tʷ nˠ aː aˑ ə̃ n̩ t̪ d̥ e̞ o̝ ɑ̟ ʊ̈ ˥˦˧˨˩}
\endgroup

% -- 5. Greek -----------------------------------------------------------------
\head{Greek}
Ἐν ἀρχῇ ἦν ὁ λόγος, καὶ ὁ λόγος ἦν πρὸς τὸν θεόν, καὶ θεὸς ἦν ὁ λόγος.
Ἄνδρα μοι ἔννεπε, Μοῦσα, πολύτροπον, ὃς μάλα πολλὰ πλάγχθη, ἐπεὶ Τροίης ἱερὸν
πτολίεθρον ἔπερσεν· πολλῶν δ’ ἀνθρώπων ἴδεν ἄστεα καὶ νόον ἔγνω.\par
\allsc{Ἡ Ὀδύσσεια τοῦ Ὁμήρου} · ᾌ ᾨ Ῥ ᾳ ῃ ῳ ϝ ϛ ϟ ϡ\par
\bi{Ἐν ἀρχῇ ἦν ὁ λόγος, καὶ ὁ λόγος ἦν πρὸς τὸν θεόν, καὶ θεὸς ἦν ὁ λόγος.}
\bi{\allsc{Ἡ Ὀδύσσεια τοῦ Ὁμήρου} · ᾌ ᾨ Ῥ ᾳ ῃ ῳ ϝ ϛ ϟ ϡ}

\head{Cyrillic}
\lng{Русский} Съешь же ещё этих мягких французских булок, да выпей чаю.\par
\bi{Съешь же ещё этих мягких французских булок, да выпей чаю.}
\bi{Всі люди народжуються вільними · Всички хора се раждат свободни · Сва људска бића · ң җ қ ҷ}

\newpage
% -- 6. features --------------------------------------------------------------
\head{Small capitals}
The \smc{unesco} report of 1954 and the \allsc{NATO} archives; \smc{Firmin
Didot} cut his types in Paris, \smc{Johann Friedrich Unger} in Berlin.
\allsc{ÉCOLE ÉLÉMENTAIRE · ZWÖLF BOXKÄMPFER · ŁÓDŹ · ÞINGVELLIR · ЖИЗНЬ}\par
\bi{The \smc{unesco} report of 1954; \smc{Firmin Didot} in Paris, \smc{Unger} in Berlin.}
\bi{\allsc{ÉCOLE ÉLÉMENTAIRE · ZWÖLF BOXKÄMPFER · ŁÓDŹ · ÞINGVELLIR · ЖИЗНЬ}}

\head{Ligatures}
Standard: fi fl ff ffi ffl fb fh fk ft — often, after, offhand, Hofburg, fifth.\par
Discretionary: {\addfontfeature{Ligatures=Rare}act, respect, Christ, histoire, aspect}.\par
\lang{German}{Historical German, with long s: Ach Gott! die Kunſt iſt lang; und
kurz iſt unſer Leben. Meſſe, ſtatt, Geſchichte, Weißſilber, Kaufhaus.}\par
\bi{Standard: fi fl ff ffi ffl fb fh fk ft — often, after, offhand, Hofburg, fifth.}
\bi{Discretionary: {\addfontfeature{Ligatures=Rare}act, respect, Christ, histoire, aspect}.}
\bi{\lang{German}{Kunſt iſt lang; Meſſe, ſtatt, Geſchichte, Weißſilber, Kaufhaus.}}

\head{Figures and symbols}
Lining 0123456789, {\osf old-style 0123456789}, fractions ½ ⅓ ¾ ⅛ ⅝,
H₂O and E = mc², x¹ x² x³, 25~°C, 12~\% and 7~‰, §~4, ¶~2, №~7, ® © ™.\par
Currencies: 1 € · £ 2 · \$ 3 · ¥ 4 · ₹ 50 · ₺ 20 · ₴ 100 · ₼ 5 · ₾ 10 · ₪ 30 · ₦ 15 · ₩ 1000 · ₿ 0.1\par
Arrows and shapes: ← → ↑ ↓ ↔ ↕ ↖ ↗ ↘ ↙ · ■ □ ▲ △ ◆ ◇ ○ ● · ⁂ ⋮ ∅\par
Mathematics: let 𝔐 be a manifold, 𝔓(X) the power set of X, 𝔖(n) the symmetric
group and 𝔭 a prime ideal; ℵ₀, ℶ₁, ℓ and 𝑙; ∑ ∏ ∫ √ ∞ ≤ ≥ ≠ ± × ÷.\par
Metrics: ⏑ ⏒ ⏓ ⏔ ⏕ ⏖ ⏗ ⏘ ⏙\par
\bi{Lining 0123456789, {\osf old-style 0123456789}, fractions ½ ⅓ ¾ ⅛ ⅝, H₂O, x² x³, §~4, №~7}
\bi{Currencies: 1 € · £ 2 · \$ 3 · ¥ 4 · ₹ 50 · ₺ 20 · ₴ 100 · ₼ 5 · ₾ 10 · ₪ 30 · ₦ 15 · ₩ 1000 · ₿ 0.1}

% -- 7. waterfall --------------------------------------------------------------
\newpage
% -- 7. the Bold ---------------------------------------------------------------
\head{Regular, Italic, Bold and Bold Italic}
\def\both#1{#1\par{\itshape #1}\par{\bfseries #1}\par{\bfseries\itshape #1}\par\medskip}
\both{It is a truth universally acknowledged, that a single man in possession
of a good fortune, must be in want of a wife.}
\both{{\Large Hamburgefonstiv HAMBURG Ærøskøbing Œuvre ẞß fi fl ff \& 1234567890}}
\both{{\Large ÁÉÍÓÚ ãẽĩõũ ǯ Ğğ Ắắ Ấấ Ǻǻ ŁłŃń Ţţ Șș}}
\both{{\Large Ἄνθρωπος ᾠδή Ῥόδος Ἤ ᾄ — Съешь же ещё мягких булок}}
\both{{\Large ɓɗƙɦʈ ʂʐɕʑ ŋəɛɔ ɹɻɺ ɡɠ Ʒʒ ʋƴ}}
\both{\allsc{Small Caps: Ærøskøbing, Đuro, Łódź} 123 {\osf 456}}
Bold in running text: the \textbf{Universal Declaration} of \textbf{1948} sets out
\textbf{rights} that apply to \textbf{every human being}, \textbf{без исключения},
\textbf{χωρίς διάκριση}, \textbf{không phân biệt}.
Italic in running text: \emph{Pride and Prejudice} (\emph{1813}), \emph{Война и мир},
\emph{Ἰλιάς}, \emph{Truyện Kiều}, \emph{Ærøskøbing}, \emph{Łódź}, \emph{Şişli}.
Bold Italic in running text: the \textbf{\emph{Dictionnaire de l’Académie}} of
\textbf{\emph{1694}}, \textbf{\emph{Война и мир}}, \textbf{\emph{Ἰλιάς}},
\textbf{\emph{Truyện Kiều}}, \textbf{\emph{Ærøskøbing}}, \textbf{\emph{Łódź}}.

\head{Sizes}
\def\sample{Harmonious typography: Ångström, Œuvre, IJ, ẞ — quixotic jumbled wizards}
{\fontsize{7}{9}\selectfont 7 pt \sample}\par
{\fontsize{9}{11}\selectfont 9 pt \sample}\par
{\fontsize{11}{14}\selectfont 11 pt \sample}\par
{\fontsize{14}{17}\selectfont 14 pt \sample}\par
{\fontsize{20}{24}\selectfont 20 pt Typography: Ångström, Œuvre}\par
{\fontsize{30}{34}\selectfont 30 pt Didot: Ἄνθρωπος}\par
{\fontsize{48}{52}\selectfont Ẫ Ợ ẞ Ж Ω \& Ǯ}

% -- 8. the Bold Italic at length, and the four styles against each other ------
\newpage
\head{Bold Italic at length}
\begingroup\osf\bfseries\itshape
When the Didot family set up their presses in Paris in the last decades of the
eighteenth century, the printed page in France had already been through two
hundred years of quiet change, and almost nobody outside the trade noticed what
happened next. François-Ambroise Didot, the father, was a printer before he was
a designer of type, and his first concern was not the shape of the letters but
the way they were measured: the point Didot, which he fixed at a little more
than the English point, gave the founders a common unit, so that a line of type
cast in one shop would stand true against a line cast in another. It sounds like
a small thing. It was not. A system of sizes that everyone could rely on meant
that a book could be set from the founts of three different foundries without
the compositor having to pack every line with slips of card, and it meant, too,
that the type itself could be made finer, because nothing had to be thickened to
survive a careless fit.

The son, Firmin, took the next step. His roman, cut in the 1780s and refined for
twenty years afterwards, pushed the contrast between thick and thin strokes
further than anyone had dared: the hairlines are almost threads, the serifs flat
and unbracketed, the stress upright, so that a capital O looks less like a
written circle than like a drawing made with a ruling pen. Readers who were used
to the warmer letters of Garamond or Caslon found it cold, and some said so in
print. Others found it exactly right for an age that admired order and clear
outlines, and the great editions of Virgil, Horace and Racine that came from the
Didot press between 1798 and 1801 were bought, admired and imitated across Europe.
\endgroup

\head{The same passage in the Italic}
\begingroup\osf\itshape
When the Didot family set up their presses in Paris in the last decades of the
eighteenth century, the printed page in France had already been through two
hundred years of quiet change, and almost nobody outside the trade noticed what
happened next. François-Ambroise Didot, the father, was a printer before he was
a designer of type, and his first concern was not the shape of the letters but
the way they were measured.
\endgroup

\newpage
\head{The four styles at 11 and 24 point}
\raggedright
\def\sampleline{Hamburgefonstiv Dictionnaire de l’Académie Zeitgeist à propos 1450}
{\fontsize{11}{14}\selectfont \sampleline\par{\itshape\sampleline}\par{\bfseries\sampleline}\par{\bfseries\itshape\sampleline}\par}
\medskip
{\fontsize{24}{29}\selectfont \sampleline\par{\itshape\sampleline}\par{\bfseries\sampleline}\par{\bfseries\itshape\sampleline}\par}

\head{The styles mixed in one line}
A roman line with an \emph{italic word}, a \textbf{\emph{bold italic word}} and a
\textbf{bold word}: the \emph{Encyclopédie}, the \textbf{\emph{Dictionnaire de
l’Académie}} and the \textbf{Code civil} of \textbf{\emph{1804}} on one shelf.

\newpage
\head{Pangrams in many languages, each in the four styles}
\newcommand\pg[3]{{\par\addfontfeature{Language=#1}\fontsize{13}{16}\selectfont\raggedright
  \textsc{\small #2}\par
  #3\par{\itshape #3}\par{\bfseries #3}\par{\bfseries\itshape #3}\par}\medskip}
\raggedright
{\small Kerning check: every line below is also a test of the kerning, pair by pair; first the pairs found unkerned or off the references.}\par\medskip
\def\kp{Ty Tw To Te Va Vo AV AT AY LT LV Kr Ra Gi Wo Yo Fa Pa 7. 4, 1. A“ “A L“ ’s »W W« T» V» “Y ‘A D” D’ AW Av Aw Ay Rv Ry RO}
{\fontsize{16}{20}\selectfont \kp\par{\itshape\kp}\par{\bfseries\kp}\par{\bfseries\itshape\kp}\par}
\bigskip
\pg{English}{English}{The quick brown fox jumps over the lazy dog.}
\pg{German}{German}{Falsches Üben von Xylophonmusik quält jeden größeren Zwerg.}
\pg{German}{German}{Zwölf Boxkämpfer jagen Viktor quer über den großen Sylter Deich.}
\pg{French}{French}{Portez ce vieux whisky au juge blond qui fume.}
\pg{French}{French}{Voix ambiguë d’un cœur qui au zéphyr préfère les jattes de kiwis.}
\pg{Spanish}{Spanish}{El veloz murciélago hindú comía feliz cardillo y kiwi.}
\pg{Spanish}{Spanish}{La cigüeña tocaba el saxofón detrás del palenque de paja.}
\pg{Portuguese}{Portuguese}{Luís argüia à Júlia que «brações, fé, chá, óxido, pôr, zângão» eram palavras do português.}
\pg{Italian}{Italian}{Quel fez sghembo copre davanti.}
\pg{Dutch}{Dutch}{Pa’s wijze lynx bezag vroom het fikse aquaduct.}
\pg{Danish}{Danish}{Quizdeltagerne spiste jordbær med fløde, mens cirkusklovnen Wolther spillede på xylofon.}
\pg{Norwegian}{Norwegian}{Vår sære Zulu fra badeøya spilte jo whist og quickstep i min taxi.}
\pg{Swedish}{Swedish}{Flygande bäckasiner söka hwila på mjuka tuvor.}
\pg{Icelandic}{Icelandic}{Kæmi ný öxi hér ykist þjófum nú bæði víl og ádrepa.}
\pg{Finnish}{Finnish}{Fahrenheit ja Celsius yrjösivät Åsan backgammon-peliin, Volkswagenissa, daiquirin ja whiskyn kera.}
\pg{Estonian}{Estonian}{Põdur Zagrebi tšellomängija-följetonist Ciqo külmetas kehvas garaažis.}
\pg{Polish}{Polish}{Pchnąć w tę łódź jeża lub ośm skrzyń fig.}
\pg{Czech}{Czech}{Příliš žluťoučký kůň úpěl ďábelské ódy.}
\pg{Slovak}{Slovak}{Kŕdeľ ďatľov učí koňa žrať kôru.}
\pg{Slovenian}{Slovenian}{Šerif bo za vajo spet kuhal domače žgance.}
\pg{Croatian}{Croatian}{Gojazni đačić s biciklom drži hmelj i finu vatu u džepu nošnje.}
\pg{Hungarian}{Hungarian}{Árvíztűrő tükörfúrógép.}
\pg{Romanian}{Romanian}{Muzicologă în bej vând whisky și tequila, preț fix.}
\pg{Turkish}{Turkish}{Pijamalı hasta yağız şoföre çabucak güvendi.}
\pg{Lithuanian}{Lithuanian}{Įlinkdama fechtuotojo špaga sublykčiojusi pragręžė apvalų arbūzą.}
\pg{Latvian}{Latvian}{Muļķa hipiji mēģina brīvi nogaršot celofāna žņaudzējčūsku.}
\pg{Irish}{Irish}{D’fhuascail Íosa, Úrmhac na hÓighe Beannaithe, pór Éava agus Ádhaimh.}
\pg{Welsh}{Welsh}{Parciais fy jac codi baw hud fy ngŵr rhag ofn i’m cŵn flino.}
\pg{Esperanto}{Esperanto}{Eĥoŝanĝo ĉiuĵaŭde.}
\pg{Vietnamese}{Vietnamese}{Tôi có thể ăn thủy tinh mà không hại gì. (a test sentence, not a strict pangram)}
\pg{Greek}{Greek}{Ξεσκεπάζω την ψυχοφθόρα βδελυγμία.}
\pg{Greek}{Greek}{Γαζίες καὶ μυρτιὲς δὲν θὰ βρῶ πιὰ στὸ χρυσαφὶ ξέφωτο.}
\pg{Russian}{Russian}{Съешь же ещё этих мягких французских булок, да выпей чаю.}
\pg{Ukrainian}{Ukrainian}{Чуєш їх, доцю, га? Кумедна ж ти, прощайся без ґольфів!}
\pg{Bulgarian}{Bulgarian}{Жълтата дюля беше щастлива, че пухът, който цъфна, замръзна като гьон.}
\pg{Serbian}{Serbian}{Љубазни фењерџија чађавог лица хоће да ми покаже штос.}
\end{document}
